\section{Further directions}
\label{sec:future}

\subsection{Beyond the section integrability threshold}

The attractive LDP is established for $0\le\gamma<1/3$ on the projective
plane, and below the corresponding normalized-potential threshold on any
positive polarized compact complex manifold. On projective space it also
holds in shrinking design tubes of every fixed order. The next objective
is a neighborhood of the canonical coupling $\gamma=1$, where the balanced
coefficient and the microscopic spectral law suggest a continuation toward
reduced Gibbs measures in complex geometry
\citep{AndreassonBermanSvensson2026}. Let $\mathcal D_1$ be the balanced
fibre on $\mathbb P^2$, and condition the reference product measure on
$\|M_1(L_N)-I_3/3\|_{\rm HS}<N^{-a}$, with fixed $0<a<1/2$.
Write this conditioned measure as $Q_N^{(1)}$, with
$N=N_k=\binom{3k+2}{2}$ as in Section~\ref{sec:equilibria}.

\begin{conjecture}[Attractive reduced Laplace principle]
\label{conj:attractive}
There is an $\epsilon>0$ such that, for
$\gamma\in(1-\epsilon,1+\epsilon)$ and every continuous
$\Phi:\mathcal P(\mathbb P^2)\to\R$, the partitions are finite for
all sufficiently large $k$ and
\begin{equation}
 \lim_{k\to\infty}\frac1N\log
 \int e^{N\gamma E_k-N\Phi(L_N)}\,dQ_N^{(1)}
 =-\inf_{\substack{\mu\in\mathcal D_1\\H(\mu)<\infty}}
       [H(\mu)-\gamma E(\mu)+\Phi(\mu)].
 \label{eq:attractive-target}
\end{equation}
\end{conjecture}

This would extend the constrained attractive LDP and its pressure to the
canonical range, including continuous empirical tilts. The first-moment coefficient
$12/7$ supplies macroscopic coercivity beyond the canonical coupling
$\gamma=1$. A microscopic route is to prove, for the attractive tilted
Gibbs laws $P_{k,\gamma,\Phi}$ and the continuous energy envelopes,
\[
 \lim_{r\to\infty}\limsup_{k\to\infty}
 \mathbb E_{P_{k,\gamma,\Phi}}
       [E_k-E_r(L_N)]_+=0.
\]
Lower comparison controls the opposite side of this residual. Through
\eqref{eq:coherent-logarithmic-energy}, the new estimate concerns
eigenvalues on exponentially small scales, rather than fixed positive
spectral moments. The directed-tree law supplies a reference-measure
benchmark; controlling its change under attractive reweighting is an
additional requirement. Extending the coupling
range toward the full macroscopic threshold is a further question.
The proof below $1/3$ identifies a specific place to seek an improvement:
compactness of the conditional densities after exploiting the global
moment constraint. The finite assignment theorem supplies an independent
computable description of torus-averaged energy, but a stronger
exponential estimate would be needed to carry it to the canonical range.

\subsection{Hermitian algebraic support and higher designs}

The cubic classification and the Hesse certificate invite a general theory
of positive design measures on determinantal varieties. A useful target
is to identify formats in which a rank witness converts some scalar
marginals into full moment-map balance, then classify the relevant closed
orbits. The equality cases $(3,3)$ and $(4,2)$ in the rank-witness
calculation show that this mechanism has a specific representation-theoretic
boundary. Beyond it, additional equations or higher moments are needed.
Quantitative ideal-to-distance bounds would turn those equations into
finite-copy ensemble certificates. Exact thresholds for higher-design
entropy fibres and the geometry of their concentrating supports are
natural problems in their own right.

The Singer construction suggests studying cyclic covers attached to
weighted designs with special vanishing divisors. Here the divisor
relation, automorphisms, and integral discriminant form determine the K3
lattice. One can ask which positive moment constraints select similarly
distinguished surfaces and whether lattice data classify the associated
ensemble faces. Such a classification would join Hermitian interpolation
to the geometry of surfaces through explicit constructions.

\subsection{Undecorated networks and spatial modular response}

The wire criterion covers arbitrary cores and uniformly primitive
attachments, including decorated AKLT models on every bounded-degree
simple graph. For undecorated two-dimensional parents, the corresponding
objective is a finite collar or regional-response estimate that implies
the required subspace angles with a uniform microscopic energy scale.
Boundary-state factorization provides a complementary setting for this
question \citep{Kastoryano2026}. The local-ratio and interaction-norm
examples show why a spatial argument must control the actual compatible
response or the support growth of modular evolution. Complete positivity
alone gives ordinary norm contraction, not an interaction-norm estimate.
This distinction is also relevant to effective Hamiltonians and strong
mixing criteria for quantum Gibbs states
\citep{BluhmCapelPerezHernandez2025}.

\subsection{Positive information algorithms and sharper contact geometry}

The clock construction admits any certified positive zero measure of the
stated form. Further input families from Lee--Yang theory could enlarge
the class of quantum score filters with positive finite approximations.
The practical questions are then quantitative: tight observable-specific
commutator bounds, optimal allocation of deterministic and sampling error,
and efficient access to the required weighted inner products. The
information-matrix bracket is a concrete target for these improvements.

The nonlinear Green formula gives a complementary computational problem
in complex geometry: obtain an effective polynomial-level bias rate and
stable interpolation rules for the random-section estimator. The
coherent-moment Hessian in \eqref{cy:hessian-hierarchy} offers positive
finite truncations for studying this question. The sampling variance is
already controlled uniformly in the level; geometric approximation error
and numerical conditioning are the additional quantities needed for a
practical metric reconstruction.

For measurements, the qualitative neighborhood at $5/12$ can be refined
by deriving explicit radii as functions of the number of outcomes and the
distance to the classified contact strata. The genuinely irreducible
qutrit region is a separate geometric problem. Combining its dual
contact inequalities with finite approximation methods for compatibility
\citep{PortoDesignollePokuttaQuintino2026} would make the theory more
directly usable in quantum-information experiments and optimization.
